\setcounter{section}{62} % tema número N-1
\section{Sistemas de gestión de Bases de Datos}

Nombre completo: Los sistemas de gestión de bases de datos SGBD. El modelo de referencia de ANSI

\localtableofcontents

\subsection{Bases de datos}
Una base de datos es un sistema capaz de trabajar con un conjunto de datos estructurados y relacionados, almacenados en un soporte informático.

Un diccionario de datos almacena información acerca de los objetos que conforman la base de datos. Su estructura, definiciones, restricciones. Es un elemento vivo donde se van incorporando elementos y metadatos. Si el diccionario está limitado a la definición existente, se llama diccionario activo. Si no se encuentran estas restricciones se denomina diccionario pasivo.


\subsection{Sistema de Gestión de Bases de Datos}
    Es un conjunto de programas, procedimientos y lenguajes que permiten almacenar, actualizar y consultar datos en una base de datos.

    \subsubsection{características y Aspectos}
        \begin{itemize}
            \item Utiliza un modelo de datos soportado por si mismo
            \item Usa lenguajes de alto nivel para interactuar y administrar
            \item Soporta acceso congruente de forma consistente y sin conflictos.
            \item Separa la parte lógica y física
            \item Redundancia controlada de datos
            \item Integridad y consistencia de datos
            \item Seguridad
            \item Alto rendimiento funcional
            \item Alto volumen de datos
            \item Recuperación (a un estado anterior íntegro)
            \item Privacidad
        \end{itemize}
        Según el teorema de CAP, en un SGBD no se pueden garantizar mas de dos de los tres aspectos:
        p=\begin{description}
            \item[Consistencia] Las lecturas son o la última escritura o un error.
            \item[Disponibilidad] Cualquier petición recibe una respuesta no erronea, pero sin garantizar la mas reciente.
            \item[Tolerancia al particionado] Seguir funcionando aún con una cantidad arbitraria de mensajes descartados entre nodos.
        \end{description}

    \subsubsection{Tipos de Lenguaje}
        \begin{description}
            \item[Data Definition Languaje DDL] Manipulación de tablas y del modelo de datos.
            \item[Data Manipulation Languaje DML] CRUD de datos
            \item[Data Control Languaje DCL] Gestión de la base de datos en sí misma.
            \item[Transaction Control Languaje TCL] Permite controlar las transacciones y concurrencia para mantener la integridad de los datos.
        \end{description}

    \subsubsection{Transacción y Concurrencia}
        Una transacción es una unidad de trabajo. En SQL se limitan entre \texttt{BEGIN-TRANSACTION} y \texttt{COMMIT-ROLBACK}.

        Estas tienen que respetar el principio ACID:
        \begin{description}
            \item[Atomicidad] La transacción se ejecuta totalmente o no se ejecuta. Commit afianza los cambios y Rollback los descarta.
            \item[Consistencia] Solo se ejecutan las operaciones que respeten la integridad de los datos
            \item[Aislamiento] Una operación no afectará a las demás.
            \item[Persistencia] Una vez se ha terminado la transacción los cambios son persistentes y sobreviven a fallos de sistema.
        \end{description}

        Se denomina concurrencia a varias operaciones pudiendo actuar al mismo tiempo. Esto genera una serie de problemas:
        \begin{description}
            \item[Lectura no repetible] Una transacción lee dos veces un valor esperando el mismo y no coincide porque una segunda transacción lo ha modificado.
            \item[Lectura Sucia] Una transacción lee un dato modificado pero no commiteado. Si la segunda transacción falla, el dato leido no es válido y por lo tanto no es consistente.
            \item[Lectura fantasma] una transacción realiza varias consultas iguales, y una segunda transacción inserta datos entre ambas consultas, dando mas datos en la segunda lectura.
        \end{description}

        Según el nivel de error aceptable se pueden establecer niveles de aislamiento.

        \begin{description}
            \item[Lectura no comprometida] Los cambios realizados se encuentran disponibles inmediatamente. No evita ningún tipo de problema
            \item[Lectura comprometida] Los cambios solo se encuentran disponibles a partir del commit. Previene las lecturas sucias.
            \item[Lectura repetible] Las filas leidas o actualizadas quedan bloqueadas hasta el fin de la transacción. Previene lecturas sucias y no repetibles.
            \item[Serializable] Las transacciones ejecutadas simultáneamente producen los mismos efectos que si se ejecutan en serie.
        \end{description}

        También existe un mecanismo llamado \textbf{Two phase locking} que obliga a ejecutar transacciones en dos fases, adquisición de recursos y liberación. De esta forma, se liberan todos de una sola vez cuando se hace commit o rollback.

    \subsubsection{Arquitectura interna}
        Se pueden distinguir tres partes. El almacenamiento en disco, el gestor de almacenamiento y el procesador de consultas. 

        El \textbf{Procesador de consultas} es la parte frontal. Analiza sintáxis y semántica de consultas. También las optimiza y las transforma en instrucciones de bajo nivel.

        El \textbf{Gestor de Transacciones} recibe múltiples transacciones concurrente y ordena las operaciones a ejecutar. También proporciona aislamiento entre transacciones.

        El \textbf{Planificador} controla concurrencia y planifica la ejecución de transacciones. Restringe el orden en el que se ejecutan las operaciones. También se denomina Gestor de Bloqueos.

        El \textbf{Gestor de Datos} Gestiona buffers de lectura y escritura. También tiene un Gestor de recuperación que maneja commit y rollback. Gestiona atomicidad, persistencia y recuperación.

\subsection{Modelo de referencia ANSI}
    Según definiciones se diferencian tres niveles.

    El invel interno o físico se encarga de las partes relacionadas con el servidor. Esta parte se dedica a mejorar el tiempo de respuesta, minimizar el espacio de almacenamiento y evitar la redundancia de datos.

    El nivel conceptual materializa la representación de datos, independientemente de su estructura física. En esta parte reside el esquema de la base de datos.

    El nivel Externo o lógico filtra los esquemas conforme a los usuarios concretos. Corresponde a la vista.

    De esta forma se garantiza la independencia de los datos y la base de datos.

    Existen tres familias de modelos de datos. El mas sencillo es el modelo jerárquico, donde se dispone en árbol y se tiene relaciones 1:n. En el modelo Codasyl hay relaciones en red. El mas extendido es el modelo relacional, donde se aplican teorías de conjuntos. Las tablas tienen relaciones con ciertas restricciones.

\subsection{Ejemplos de bases de datos}

\begin{itemize}
    \item Relacionales
    \begin{itemize}
        \item Oracle
        \item SQL-Server
        \item MySQL
        \item MariaDB
        \item DB2
        \item PostgreSQL
    \end{itemize}
    \item Orientadas a objetos
    \begin{itemize}
        \item ObjectDB
        \item Zope
        \item Perst
        \item ObjectStore
        \item DB4O
    \end{itemize}
    \item Orientadas a Grafos
    \begin{itemize}
        \item Allegro
        \item Neo4j
        \item Sparksee
        \item ArangoDB
        \item OrientDB
    \end{itemize}
    \item No-SQL
        \begin{itemize}
            \item Clave valor
            \begin{itemize}
                \item DynamoDB
                \item ArangoDB
                \item OrientDB
                \item Cassandra
                \item Memcached
                \item Redis
            \end{itemize}
            \item Orientadas a Columnas
            \begin{itemize}
                \item Cassandra
                \item HBase
            \end{itemize}
            \item Documentales
            \begin{itemize}
                \item MongoDB
                \item BaseX
                \item MarkLogic
                \item CouchDB
                \item RavenDB
                \item Arango
                \item OrientDB
                \item eXistdb
            \end{itemize}
        \end{itemize}
    \item Multidimensionales
\end{itemize}